\section[Pembuktian dengan Kasus]{\raggedright Pembuktian dengan Kasus dan Pembuktian Exhaustif}

Terdapat situasi di mana membuktikan suatu pernyataan secara umum (untuk seluruh elemen domain) lebih mudah dilakukan dengan membagi domain menjadi beberapa kasus yang saling lepas (\textit{mutually exclusive}) dan menangkap seluruh kemungkinan (\textit{collectively exhaustive}), kemudian membuktikan pernyataan tersebut untuk setiap kasus secara terpisah \cite{rosen2019}.

\subsection{Pembuktian dengan Kasus (Proof by Cases)}

\begin{definisi}[Pembuktian dengan Kasus]
\logic{Pembuktian dengan kasus} (\textit{proof by cases} / \textit{proof by exhaustion of cases}) adalah teknik pembuktian di mana domain dipartisi menjadi kasus-kasus $C_1, C_2, \ldots, C_n$ yang saling lepas dan lengkap, kemudian pernyataan dibuktikan secara terpisah untuk setiap kasus. Teknik ini valid berdasarkan aturan inferensi:
\[ \frac{(C_1 \rightarrow q) \wedge (C_2 \rightarrow q) \wedge \ldots \wedge (C_n \rightarrow q)}{(C_1 \vee C_2 \vee \ldots \vee C_n) \rightarrow q} \]
\end{definisi}

Kunci keberhasilan metode ini adalah memastikan bahwa kasus-kasus yang dipilih benar-benar \emph{mencakup seluruh kemungkinan}. Jika ada kasus yang terlewat, bukti menjadi tidak valid.

\begin{contoh}
\textbf{Teorema}: Untuk setiap bilangan bulat $n$, $n^2 + n$ adalah bilangan genap.

\textbf{Bukti (dengan Kasus):}

Setiap bilangan bulat $n$ bernilai genap atau ganjil (dua kasus, saling lepas dan lengkap).

\textbf{Kasus 1:} $n$ genap. Maka $n = 2k$ untuk suatu bilangan bulat $k$.
\[ n^2 + n = (2k)^2 + 2k = 4k^2 + 2k = 2(2k^2 + k) \]
yang berbentuk $2m$ (genap). $\checkmark$

\textbf{Kasus 2:} $n$ ganjil. Maka $n = 2k + 1$ untuk suatu bilangan bulat $k$.
\[ n^2 + n = (2k+1)^2 + (2k+1) = 4k^2 + 4k + 1 + 2k + 1 = 4k^2 + 6k + 2 = 2(2k^2 + 3k + 1) \]
yang berbentuk $2m$ (genap). $\checkmark$

Kedua kasus terbukti. $\blacksquare$
\end{contoh}

\subsection{Pembuktian Exhaustif}

\begin{definisi}[Pembuktian Exhaustif]
\logic{Pembuktian exhaustif} (\textit{exhaustive proof} / \textit{brute-force proof}) adalah pembuktian di mana pernyataan diverifikasi untuk \textbf{setiap} elemen domain secara individual, tanpa menggunakan argumentasi umum. Metode ini hanya dapat diterapkan jika domain bersifat berhingga dan cukup kecil untuk diperiksa satu per satu.
\end{definisi}

\begin{contoh}
\textbf{Teorema}: Untuk setiap bilangan bulat $n$ dengan $1 \leq n \leq 4$, berlaku $n^2 \leq 2^n \cdot n$.

\textbf{Bukti (Exhaustif):} Periksa setiap nilai:
\begin{center}
\renewcommand{\arraystretch}{1.2}
\begin{tabular}{|c|c|c|c|}
\hline
$n$ & $n^2$ & $2^n \cdot n$ & $n^2 \leq 2^n \cdot n$? \\ \hline
1 & 1 & 2 & Ya ($1 \leq 2$) \\ \hline
2 & 4 & 8 & Ya ($4 \leq 8$) \\ \hline
3 & 9 & 24 & Ya ($9 \leq 24$) \\ \hline
4 & 16 & 64 & Ya ($16 \leq 64$) \\ \hline
\end{tabular}
\end{center}

Semua kasus terverifikasi. $\blacksquare$
\end{contoh}

\begin{catatan}
Pembuktian exhaustif memiliki keterbatasan yang jelas: jika domain berhingga tetapi sangat besar (misalnya semua bilangan prima kurang dari $10^9$), verifikasi manual menjadi tidak praktis. Dalam kasus seperti itu, komputer sering digunakan untuk membantu verifikasi, tetapi bukti semacam itu kadang diperdebatkan statusnya dalam komunitas matematika.
\end{catatan}
